\scanpage{272}\mspage{188}
% archive sheet number: 271
c') Comme c), \deleted{\unsure{mais}} avec $\pi_0$ cardinal
régulier (et non seulement utile) pour la
filtration card. gén. \inserted{acc.} en question, et cette
\deleted{\unsure{la}} filtration étant supposée de plus \uline{canonique}.
% margin note written sideways
\margin{On a alors le diagramme d'implications
\diagram{\begin{tikzcd}[ampersand replacement=\&,column sep=3pt,cells={nodes={inner sep=1pt,font=\scriptsize}}]
c'\arrow[r,Leftrightarrow]\&c\arrow[r,Leftrightarrow]\&b\arrow[r,Rightarrow]\&a
\end{tikzcd}}}

4.8.1.1. Il me faut ici préciser la
terminologie, sur ce que j'entends par
sous-catégorie (pleine) $C$ \uline{strictement} \uline{génératrice}
dans $M$. Pour ceci, je me réfère à
SGA 4 I 7.2, (i), aux conditions équivalentes b) et c) \unsure{du|de} loc. cit.

\textcircled{b} $\forall F\in\operatorname{Ob}M$, la flèche canonique
\[
\varinjlim_{i\in C/F}x_i\longrightarrow F
\]
est un iso ;

\textcircled{c} Le foncteur canonique
\[
\begin{aligned}
M&\longrightarrow C^\wedge\\
F&\longmapsto\bigl(\underset{\substack{\unsure{\rotatebox{90}{\(\in\)}}\\\operatorname{Ob}C}}{x}
\unsure{\mapstochar\relbar|\mapsto}\underset{\in\operatorname{Ob}(\mathrm{Ens})}{\underline{\operatorname{Hom}_M(x,F)}}\bigr).
\end{aligned}
\]
est pleinement fidèle.

Ces conditions impliquent

a) $C$ est génératrice \inserted{dans $M$} par épimorphismes stricts, \inserted{i.e. pour tt $F$ \unsure{dans|de} $M$, la famille des flèches $x\to F$ de source dans $C$, de but $F$, est épimorphique stricte.}
et lorsque $M$ est stable par produits
fibrés, cette condition a) est équivalente
aux précédentes. Dans loc. cit. j'affirme
